\documentclass[11pt]{article}
\usepackage[a4paper,margin=1in]{geometry}
\usepackage{fontspec}
\usepackage[boldfont=false]{xeCJK}
\setCJKmainfont{FandolSong-Regular.otf}
\usepackage{amsmath}
\usepackage{amssymb}
\usepackage{array}
\usepackage{booktabs}
\usepackage{graphicx}
\usepackage{adjustbox}
\usepackage{enumitem}
\setlist[itemize]{topsep=2pt,partopsep=0pt,parsep=0pt,itemsep=2pt,leftmargin=1.4em}
\setlist[enumerate]{topsep=2pt,partopsep=0pt,parsep=0pt,itemsep=2pt,leftmargin=1.6em}
\usepackage{parskip}
\usepackage{fvextra}
\fvset{breaklines=true,breakanywhere=true,fontsize=\small}
\setlength{\parindent}{0pt}

\begin{document}

\begin{center}
  {\LARGE\bfseries Atoms, elements and compounds}\\[0.35em]
  {\large Igcse Chemistry}
\end{center}
\vspace{0.6em}

\section*{Elements, compounds and mixtures}

All substances are made from about 100 simple building blocks. Knowing how they are joined lets you sort every substance into one of three groups.

\begin{itemize}
\item An \textbf{element} 元素 is a substance made of only one type of \textbf{atom} 原子. You cannot break it into anything simpler by a chemical reaction. Examples: copper, oxygen, carbon.

\item A \textbf{compound} 化合物 is two or more elements chemically joined (bonded) together. The atoms are joined in a fixed ratio. Examples: water, carbon dioxide. A compound has different properties from the elements in it.

\item A \textbf{mixture} 混合物 is two or more substances that are just mixed, not chemically joined. The parts keep their own properties and can be separated by physical methods. Example: air.

\end{itemize}

The key difference: in a compound the elements are \textbf{bonded} and can only be separated by chemical reactions; in a mixture they are not bonded and are easy to separate.

\section*{Atomic structure}

\subsection*{Inside the atom}

Every atom has a small, dense centre called the \textbf{nucleus} 原子核. Around it, \textbf{electrons} 电子 move in \textbf{shells} 电子层 (energy levels). The nucleus contains two kinds of particle: \textbf{protons} 质子 and \textbf{neutrons} 中子.

Each particle has a \textbf{relative mass} and a \textbf{relative charge} 电荷. You must learn these values:

\begin{center}
\adjustbox{max width=\linewidth}{%
\begin{tabular}{lll}
\toprule
\textbf{Particle} & \textbf{Relative mass} & \textbf{Relative charge} \\
\midrule
proton & 1 & $+1$ \\
neutron & 1 & $0$ \\
electron & $\frac{1}{1840}$ (almost 0) & $-1$ \\
\bottomrule
\end{tabular}}
\end{center}

An atom has no overall charge because it has equal numbers of protons ($+1$ each) and electrons ($-1$ each).

\subsection*{Proton number and mass number}

Two numbers describe an atom:

\begin{itemize}
\item The \textbf{proton number} 质子数 (also called the \textbf{atomic number} 原子序数) is the number of protons in the nucleus. It tells you which element the atom is.

\item The \textbf{mass number} 质量数 (also called the \textbf{nucleon number} 核子数) is the total number of protons and neutrons in the nucleus.

\end{itemize}

So the number of neutrons $=$ mass number $-$ proton number.

\subsection*{Electronic configuration}

The electrons fill the shells from the inside out. The first shell holds up to 2 electrons; the next shells hold up to 8 each (for the first 20 elements). The \textbf{electronic configuration} 电子排布 lists how many electrons are in each shell, starting from the inside.

For example, an atom with 13 electrons has the configuration $2,8,3$. Sodium (proton number 11) is $2,8,1$. Calcium (proton number 20) is $2,8,8,2$.

The configuration links to the \textbf{Periodic Table} 周期表:

\begin{itemize}
\item A \textbf{Group} 族 number (Groups I to VII) equals the number of electrons in the outer shell. So $2,8,1$ is in Group I.

\item A \textbf{Period} 周期 number equals the number of shells that hold electrons. So $2,8,1$ has three shells, so it is in Period 3.

\item The \textbf{noble gases} 稀有气体 in Group VIII (or 0) have a full outer shell, which makes them very unreactive.

\end{itemize}

\section*{Isotopes}

\textbf{Isotopes} 同位素 are atoms of the same element that have the same number of protons but different numbers of neutrons. Because the proton number is the same, they are the same element. Because the neutron number is different, they have different mass numbers.

You write an atom with its mass number on top and proton number below, like $^{12}_{6}\text{C}$. For an ion you add the charge, like $^{35}_{17}\text{Cl}^{-}$.

Isotopes of an element have the \textbf{same chemical properties}. This is because chemical reactions only involve electrons, and isotopes have the same number of electrons and the same electronic configuration. (The extra neutrons change only the mass.)

\subsection*{Calculating relative atomic mass}

The \textbf{relative atomic mass} 相对原子质量 ($A_r$) of an element is the average mass of its atoms, taking into account how common each isotope is. The \textbf{abundance} 丰度 is the percentage of each isotope.

\[
A_r = \frac{\sum (\text{isotope mass} \times \text{abundance})}{100}
\]

For example, chlorine is 75\% $^{35}\text{Cl}$ and 25\% $^{37}\text{Cl}$:

\[
A_r = \frac{(35 \times 75) + (37 \times 25)}{100} = 35.5
\]

\section*{Ions and ionic bonds}

An \textbf{ion} 离子 is an atom (or group of atoms) that has lost or gained electrons, so it has an electric charge.

\begin{itemize}
\item A metal atom \textbf{loses} electrons to form a positive ion, called a \textbf{cation} 阳离子.

\item A non-metal atom \textbf{gains} electrons to form a negative ion, called an \textbf{anion} 阴离子.

\end{itemize}

Atoms do this to get a full outer shell, like a noble gas.

\subsection*{How an ionic bond forms}

An \textbf{ionic bond} 离子键 is a strong \textbf{electrostatic attraction} 静电引力 between oppositely charged ions (a $+$ ion and a $-$ ion pull together).

Ionic bonds form between a \textbf{metal} 金属 and a \textbf{non-metal} 非金属. Take sodium chloride, $\text{NaCl}$. Sodium ($2,8,1$) gives its one outer electron to chlorine ($2,8,7$). Now sodium is $\text{Na}^{+}$ ($2,8$) and chlorine is $\text{Cl}^{-}$ ($2,8,8$). Both have full outer shells, and the opposite charges attract.

You can show this with a \textbf{dot-and-cross diagram}: draw each atom's outer-shell electrons as dots for one element and crosses for the other, then show the electron moving from the metal to the non-metal.

\subsection*{The structure and properties of ionic compounds}

An \textbf{ionic compound} 离子化合物 is not made of separate \textbf{molecules} 分子. The ions pack together into a giant \textbf{lattice} 晶格 — a regular pattern of huge numbers of \textbf{alternating} 交替 positive and negative ions.

This structure explains the properties:

\begin{center}
\adjustbox{max width=\linewidth}{%
\begin{tabular}{ll}
\toprule
\textbf{Property} & \textbf{Reason} \\
\midrule
high \textbf{melting point} 熔点 and \textbf{boiling point} 沸点 & the strong electrostatic attraction between ions needs a lot of energy to break \\
poor \textbf{electrical conductivity} 导电性 when solid & the ions are fixed in place and cannot move \\
good conductor when \textbf{molten} 熔融 or \textbf{aqueous} 水溶液 & the ions are now free to move and carry charge \\
\bottomrule
\end{tabular}}
\end{center}

\section*{Simple molecules and covalent bonds}

A \textbf{covalent bond} 共价键 forms when two atoms \textbf{share} a pair of electrons. By sharing, each atom gets a full outer shell (a noble gas configuration). Covalent bonds form between non-metal atoms.

Some molecules to know:

\begin{itemize}
\item $\text{H}_2$ — two hydrogen atoms share one pair of electrons (a single bond).

\item $\text{Cl}_2$, $\text{HCl}$ — one shared pair each.

\item $\text{H}_2\text{O}$ — oxygen shares one pair with each of two hydrogen atoms.

\item $\text{NH}_3$ — nitrogen shares a pair with each of three hydrogen atoms.

\item $\text{CH}_4$ — carbon shares a pair with each of four hydrogen atoms.

\item $\text{O}_2$ and $\text{CO}_2$ have double bonds (two shared pairs); $\text{N}_2$ has a triple bond (three shared pairs); $\text{C}_2\text{H}_4$ and $\text{CH}_3\text{OH}$ also use shared pairs.

\end{itemize}

In a dot-and-cross diagram for a molecule, you draw the outer electrons of each atom and show which pairs are shared in the overlap between the atoms.

\subsection*{Properties of simple molecular compounds}

These substances are made of small, separate molecules.

\begin{itemize}
\item They have \textbf{low} melting points and boiling points. The covalent bonds inside each molecule are strong, but the \textbf{intermolecular forces} 分子间作用力 (the forces between one molecule and the next) are weak, so little energy is needed to separate the molecules.

\item They are \textbf{poor conductors} of electricity, because the molecules have no overall charge and no free electrons or ions to carry charge.

\end{itemize}

\section*{Giant covalent structures}

Some covalent substances are not small molecules. Instead, millions of atoms are joined by covalent bonds into one \textbf{giant covalent structure} 巨型共价结构. The two you must know are both forms of carbon.

\textbf{Diamond} 金刚石: each carbon atom is bonded to four other carbon atoms. This makes a very strong, rigid 3-D network. Diamond is extremely hard, so it is used in \textbf{cutting tools} 切割工具.

\textbf{Graphite} 石墨: each carbon atom is bonded to only three others, forming flat \textbf{layers} 层. There are weak forces between the layers, so the layers can slide over each other — this makes graphite a good \textbf{lubricant} 润滑剂. The fourth outer electron of each carbon is free; these \textbf{delocalised electrons} 离域电子 can move, so graphite conducts electricity and is used as an \textbf{electrode} 电极.

\textbf{Silicon(IV) oxide} 二氧化硅 ($\text{SiO}_2$) has a giant covalent structure like diamond, so it is also very hard and has a very high melting point.

\section*{Metallic bonding}

A metal is a giant structure of positive ions surrounded by a 'sea' of \textbf{delocalised electrons} that are free to move through the whole metal. \textbf{Metallic bonding} 金属键 is the strong electrostatic attraction between these positive ions and the sea of electrons.

This explains two key properties of metals:

\begin{itemize}
\item \textbf{Good electrical conductivity}: the delocalised electrons are free to move and carry charge through the metal.

\item \textbf{Malleability} 展性 (can be hammered into sheets) and \textbf{ductility} 延性 (can be pulled into wires): the layers of positive ions can slide over each other without breaking the metallic bond, so the metal changes shape instead of shattering.

\end{itemize}


\end{document}
